\documentclass[11pt,a4paper]{article}
\usepackage[top=1.7cm,bottom=1.5cm,left=1.5cm,right=1.5cm,headsep=8pt,footskip=16pt]{geometry}
\usepackage{enumitem}
\usepackage{amssymb}
\usepackage{array}
\usepackage[table]{xcolor}
\usepackage{fancyhdr}
\usepackage{lastpage}
\usepackage{multicol}
\usepackage{polyglossia}
\setdefaultlanguage[calendar=gregorian, numerals=maghrib]{arabic}
\setotherlanguage{english}
\IfFontExistsTF{Amiri}{%
  \newfontfamily\arabicfont[Script=Arabic, Scale=1.08]{Amiri}
}{%
  \newfontfamily\arabicfont[Script=Arabic, Scale=1.08, Path=../fonts/, Extension=.ttf,
    UprightFont=*-Regular, BoldFont=*-Bold, ItalicFont=*-Italic,
    BoldItalicFont=*-BoldItalic]{Amiri}
}
\IfFontExistsTF{AlBattar}{%
  \newfontfamily\khatt[Script=Arabic, Scale=1.15]{AlBattar}
}{%
  \newfontfamily\khatt[Script=Arabic, Scale=1.15, Path=../fonts/,
    BoldFont=ae_AlBattar.ttf]{ae_AlBattar.ttf}
}
\renewcommand{\headrulewidth}{0.8pt}
\renewcommand{\footrulewidth}{0.8pt}
\fancyhf{}
\rhead{\khatt المدرسة العليا للأساتذة ببوزريعة}
\lhead{\khatt ملحقة بومرداس}
\lfoot{\khatt الأستاذ: د. بن عبد الرزاق يوسف}
\cfoot{\thepage/\pageref{LastPage}}
\pagestyle{fancy}
\parindent 0pt
\setlength{\columnsep}{0.6cm}
\setlength{\columnseprule}{0.4pt}
\setlist[enumerate,1]{label=\textbf{\arabic*.}, leftmargin=*, itemsep=5pt, topsep=2pt}
\begin{document}

\linespread{1.08}\selectfont
\begin{center}
{\LARGE\khatt\textbf{الاختبار الحادي عشر في الذكاء الاصطناعي}}\\[4pt]
{\normalsize وضعيات إشكالية ــ أسئلة متعددة الاختيارات ــ الفصول من 1 إلى 6}
\end{center}
\vspace{-2pt}\rule{\textwidth}{0.6pt}\\[2pt]
\textbf{التعليمة:} يتضمن الاختبار خمس وضعيات إشكالية. اقرأ كل وضعية بعناية، ثم اختر التحليل أو التصرف الأنسب
من بين الاختيارات الخمسة؛ لكل وضعية إجابة صحيحة واحدة فقط.
اكتب \textbf{حرف} الإجابة الصحيحة في جدول الإجابات في آخر الاختبار.
لكل وضعية 4 نقاط (المجموع 20 نقطة).
\vspace{4pt}

\noindent
{\large\khatt\textbf{الوضعية 1}}\hfill{\small(الفصل 1)}\par\nopagebreak\vspace{3pt}
\setlength{\fboxsep}{7pt}\colorbox{gray!10}{\parbox{\dimexpr\linewidth-14pt}{يستعمل أستاذ برنامجًا لتدقيق الإملاء منذ سنوات، ثم جرّب أداة محادثة توليدية فوجدها تعيد كتابة الفقرة كلها بأسلوب جديد، بينما كان البرنامج القديم يكتفي بتسطير الأخطاء. ويريد أن يشرح الفرق لزملائه في يوم دراسي. ويتساءل: هل الأداتان من النوع نفسه؟ وأيهما أنسب لتقويم إنشاء التلميذ دون المساس بأسلوبه؟}}\par\nopagebreak
\begin{itemize}[leftmargin=1.0cm, labelsep=6pt, itemsep=3pt, topsep=4pt, parsep=0pt]
  \item[\textbf{أ)}] الأداتان مختلفتان، والأنسب للتقويم إعادة كتابة الإنشاء بالأداة التوليدية ثم منح التلميذ علامة النص المعاد كتابته، لأن العلامة يجب أن تعكس أفضل صيغة ممكنة لأفكاره، فينال التلميذ ما يستحقه من تشجيع.
  \item[\textbf{ب)}] البرنامج القديم يرصد الأخطاء وفق قواعد أو نماذج محددة، والأداة التوليدية تنتج نصًا جديدًا؛ ولتقويم إنشاء التلميذ يُستعمل المدقق، أو تُطلب من الأداة ملاحظات دون إعادة كتابة، حفاظًا على أسلوبه.
  \item[\textbf{ج)}] الأداة التوليدية ترصد الأخطاء فقط مثل البرنامج القديم، وإعادة الكتابة التي لاحظها الأستاذ خلل مؤقت، فتُستعمل لتقويم الإنشاءات كما هي، دون تحديد ما يريده في الطلب.
  \item[\textbf{د)}] الأداتان من النوع نفسه تمامًا، والفرق الوحيد في السرعة، ولذلك تُستعمل الأداة التوليدية لإعادة كتابة كل الإنشاءات، لأن النص المعاد كتابته أفضل دائمًا من نص التلميذ الأصلي، ويوفر على الأستاذ جهد التصحيح كله.
  \item[\textbf{هـ)}] البرنامج القديم ذكاء اصطناعي توليدي أيضًا، لأنه يولّد خطوطًا حمراء تحت الكلمات، فلا فرق بين الأداتين في المبدأ، ويمكن استعمال أي منهما لإعادة كتابة إنشاء التلميذ.
\end{itemize}
\par\vspace{6pt}
\noindent
{\large\khatt\textbf{الوضعية 2}}\hfill{\small(الفصل 2)}\par\nopagebreak\vspace{3pt}
\setlength{\fboxsep}{7pt}\colorbox{gray!10}{\parbox{\dimexpr\linewidth-14pt}{طلب أستاذ ترجمة جملة إنجليزية تعني «الطبيب تحدث مع الممرضة» إلى العربية في سياق محايد، فترجمتها الأداة مذكّرًا للطبيب ومؤنثًا للممرضة دائمًا، مع أن الأصل الإنجليزي لا يحدد الجنس. وتكررت الظاهرة مع «المهندس» و«المعلمة». ويريد استعمال الترجمات في نصوص قرائية يوزعها على تلاميذه. كيف تُفسَّر هذه الظاهرة؟}}\par\nopagebreak
\begin{itemize}[leftmargin=1.0cm, labelsep=6pt, itemsep=3pt, topsep=4pt, parsep=0pt]
  \item[\textbf{أ)}] العربية تفرض أن يكون الطبيب مذكرًا والممرضة مؤنثة، ولذلك فالترجمة هي الوحيدة الصحيحة نحويًا، ولا علاقة للظاهرة ببيانات التدريب أو بأي انحياز في النموذج.
  \item[\textbf{ب)}] النموذج مبرمج يدويًا بقائمة تربط كل مهنة بجنس محدد وضعها المطوّرون عمدًا، فلا يتعلم شيئًا من النصوص، والحل الوحيد مراسلة الشركة لتعديل هذه القائمة الثابتة.
  \item[\textbf{ج)}] النموذج تعلّم من نصوص يغلب فيها اقتران مهن بجنس معين، فيعيد إنتاج هذا النمط عند الغموض؛ وهو انحياز موروث من البيانات، يعالجه المستخدم بتحديد الجنس في الطلب أو بمراجعة الترجمة.
  \item[\textbf{د)}] المترجم الآلي يختار الجنس عشوائيًا في كل مرة، وتكرار النمط نفسه مصادفة لا تدل على شيء، ولو أعاد الأستاذ الترجمة مرات كافية لتوزع الجنسان بالتساوي على المهن كلها، فالأفضل الاكتفاء بالترجمة الأولى.
  \item[\textbf{هـ)}] النموذج تعلّم من نصوص يغلب فيها هذا الاقتران، وهذا يدل على أن الواقع كذلك بالضرورة، فتُعتمد الترجمة كما هي دون مراجعة، لأن البيانات تعكس الحقيقة الاجتماعية بدقة، ولا حاجة إلى تعديلها في النصوص القرائية.
\end{itemize}
\par\vspace{6pt}
\noindent
{\large\khatt\textbf{الوضعية 3}}\hfill{\small(الفصل 3)}\par\nopagebreak\vspace{3pt}
\setlength{\fboxsep}{7pt}\colorbox{gray!10}{\parbox{\dimexpr\linewidth-14pt}{يريد أستاذ أن تساعده الأداة في إعراب جمل معقدة لطلبته، فلاحظ أنها حين تعطي الإعراب مباشرة تخطئ أحيانًا في تحديد العامل. واقترح زميله أن يغيّر طريقة الطلب. والأستاذ يريد إعرابًا يمكن تتبّعه والتحقق منه. ويحتاج إلى نماذج إعراب يعرضها على طلبته ويشرحها خطوة خطوة. أي صياغة للطلب أنسب لهذا الغرض؟}}\par\nopagebreak
\begin{itemize}[leftmargin=1.0cm, labelsep=6pt, itemsep=3pt, topsep=4pt, parsep=0pt]
  \item[\textbf{أ)}] «أعرب الجملة خطوة خطوة وعلّل كل حكم»، ثم يعتمد الإعراب دون مراجعة، لأن وجود التعليل دليل كافٍ على صحة الإعراب، فالأداة لا تعلّل حكمًا خاطئًا أبدًا، وهذا يوفر وقت الأستاذ وجهده في المراجعة.
  \item[\textbf{ب)}] «أعرب الجملة في كلمة واحدة لكل لفظ دون تعليل»، لأن الإجابة المختصرة أقل عرضة للخطأ من الإجابة المفصلة، والتعليل يفتح للأداة بابًا لمزيد من الأخطاء النحوية.
  \item[\textbf{ج)}] «أعرب الجملة إعرابًا صحيحًا ولا تخطئ»، لأن النهي الصريح عن الخطأ يجعل الأداة أكثر حرصًا، ويغني عن تفصيل خطوات الإعراب أو تعليل الأحكام في الإجابة.
  \item[\textbf{د)}] «أعرب الجملة كما يعربها سيبويه»، لأن ذكر اسم نحوي كبير يجعل الأداة تسترجع إعرابه الحرفي للجملة نفسها من كتابه، فيأتي الإعراب صحيحًا دون حاجة إلى تفصيل أو مراجعة، كما هو في كتب الإعراب المشهورة.
  \item[\textbf{هـ)}] «أعرب الجملة خطوة خطوة: حدّد أولًا نوعها، ثم الفعل أو المبتدأ، ثم العامل في كل كلمة، وعلّل كل حكم إعرابي»، ثم يراجع الخطوات بنفسه، لأن تفصيل الاستدلال يكشف موضع الخطأ إن وقع.
\end{itemize}
\par\vspace{6pt}
\noindent
{\large\khatt\textbf{الوضعية 4}}\hfill{\small(الفصل 4)}\par\nopagebreak\vspace{3pt}
\setlength{\fboxsep}{7pt}\colorbox{gray!10}{\parbox{\dimexpr\linewidth-14pt}{استعمل أستاذ برنامجًا يدّعي كشف النصوص المولَّدة بالذكاء الاصطناعي، فأظهر أن إنشاء طالبة «مولَّد بنسبة 92\%». أنكرت الطالبة، وقدّمت مسودات بخط يدها تُظهر مراحل كتابة النص. ويريد الأستاذ اتخاذ قرار عادل قبل اجتماع اللجنة. ويطالب بعض أعضاء اللجنة بالاكتفاء بنتيجة البرنامج لأنها «علمية». ما الموقف الأسلم؟}}\par\nopagebreak
\begin{itemize}[leftmargin=1.0cm, labelsep=6pt, itemsep=3pt, topsep=4pt, parsep=0pt]
  \item[\textbf{أ)}] لا تكفي نتيجة الكاشف وحدها لإدانة الطالبة، لأن هذه البرامج تخطئ كثيرًا، خاصة مع النصوص العربية والأساليب المدرسية؛ فتُعتمد قرائن أخرى كالمسودات والمحاورة الشفهية حول النص.
  \item[\textbf{ب)}] لا تُدان الطالبة، لكن يُلغى استعمال الإنشاءات المنزلية نهائيًا في التقويم، لأن وجود الكاشفات يجعل كل إنشاء منزلي موضع شك، ولا يمكن الوثوق بأي قرينة أخرى.
  \item[\textbf{ج)}] يُستعمل كاشف ثانٍ، فإن أعطى نسبة مرتفعة أيضًا ثبت الغش، لأن اتفاق كاشفين على النتيجة نفسها دليل قاطع، ولا قيمة للمسودات الورقية أمام التحليل الآلي.
  \item[\textbf{د)}] تُعتمد نتيجة الكاشف لأنها رقمية ودقيقة، والمسودات قد تكون منسوخة لاحقًا من النص المولَّد، فنسبة 92\% دليل قاطع يكفي لإلغاء الإنشاء دون حاجة إلى محاورة الطالبة.
  \item[\textbf{هـ)}] تُخفَّض علامة الطالبة بنسبة 92\% من العلامة الكاملة، لأن الكاشف حدّد مقدار الاستعمال بدقة، فيكون التخفيض متناسبًا مع حجم المخالفة دون حاجة إلى قرائن إضافية، فيكون الجزاء عادلًا ومتناسبًا.
\end{itemize}
\par\vspace{6pt}
\noindent
{\large\khatt\textbf{الوضعية 5}}\hfill{\small(الفصل 5)}\par\nopagebreak\vspace{3pt}
\setlength{\fboxsep}{7pt}\colorbox{gray!10}{\parbox{\dimexpr\linewidth-14pt}{طلب أستاذ من أداة تلخيص كتاب من ثلاثمئة صفحة رفعه ملفًا واحدًا، فجاء الملخص دقيقًا في الفصول الأولى، لكنه عام وغامض في الفصول الأخيرة، وأغفل الخاتمة تقريبًا. ويحتاج ملخصًا لكل فصل يعرضه على طلبته. ولا يملك وقتًا لقراءة الكتاب كاملًا من جديد، ويريد ملخصات أمينة يعتمد عليها في التدريس. ما التفسير والإجراء الأنسب؟}}\par\nopagebreak
\begin{itemize}[leftmargin=1.0cm, labelsep=6pt, itemsep=3pt, topsep=4pt, parsep=0pt]
  \item[\textbf{أ)}] الأداة لخّصت الفصول الأخيرة باختصار لأنها أقل أهمية دائمًا، فالمؤلفون يضعون أهم أفكارهم في البداية، ولذلك يُعتمد الملخص كما هو دون إعادة أو مقابلة.
  \item[\textbf{ب)}] يعيد الطلب نفسه بالملف نفسه مضيفًا «لخّص كل الفصول بالتساوي»، لأن هذه العبارة كافية لتجاوز أي حدود تقنية، ويعتمد الملخص الجديد دون مقابلته بالكتاب.
  \item[\textbf{ج)}] يلخّص كل فصل في طلب مستقل، ثم يجمع الملخصات ويعرضها مباشرة دون مقابلتها بالفصول، لأن تقسيم الكتاب وحده يضمن أمانة كل ملخص في كل الأحوال، فيوفر بذلك وقتًا كبيرًا قبل الحصة.
  \item[\textbf{د)}] قد يتجاوز الكتاب حدود ما تعالجه الأداة بدقة في طلب واحد، فيضعف التركيز على آخره؛ والأنسب تلخيص كل فصل في طلب مستقل، ثم مقابلة كل ملخص بفصله قبل جمعها.
  \item[\textbf{هـ)}] المشكلة في صيغة الملف، فيحوّله إلى صور ويرفعه من جديد، لأن الأداة تقرأ الصور أدق من النصوص، فيأتي الملخص متوازنًا بين الفصول دون حاجة إلى تقسيم الكتاب، وهي طريقة معروفة عند كثير من المستخدمين.
\end{itemize}
\par\vspace{6pt}
\begin{center}
{\khatt\textbf{جدول الإجابات}}\\[3pt]
\renewcommand{\arraystretch}{1.6}
\begin{tabular}{|>{\centering\arraybackslash}m{1.9cm}|*{5}{>{\centering\arraybackslash}m{1.2cm}|}}
\hline
\textbf{رقم السؤال} & \textbf{1} & \textbf{2} & \textbf{3} & \textbf{4} & \textbf{5} \\ \hline
\textbf{الإجابة} &  &  &  &  &  \\[8pt] \hline
\end{tabular}

\vspace{8pt}
\renewcommand{\arraystretch}{2.2}
\begin{tabular}{|>{\centering\arraybackslash}m{7.5cm}|>{\centering\arraybackslash}m{3.2cm}|>{\centering\arraybackslash}m{3.2cm}|}
\hline
\textbf{الاسم واللقب} & \textbf{الفوج} & \textbf{العلامة} \\ \hline
 & & \hfill /20 \\ \hline
\end{tabular}
\end{center}
\end{document}
